% 3-7-2-lstm.tex
%  Budget: 4pp.  Converted from thesis/drafts/3-7-2-lstm.md.
%  Source map: recurrent architectures for temporal dependence in EPF
%   -> zhang_deep_2020 , ugurlu_electricity_2018 ; deep models as the next
%   generation after classical ML -> oconnor_review_2025 .
%  Hyperparameters verified 2026-09-05 against configs/tuned/lstm_params.yaml:
%   units 29, epochs 47, batch_size 32, lr 0.00083171, validation_mae 3.696.
%  Test numbers from results_canonical.csv: MAE 3.873, rMAE 0.424.
%  OOD from ood_stress.csv: LSTM rMAE 1.520 vs LEAR-LASSO 1.087.
%
%  *** CORRECTED FROM THE DRAFT ***  The markdown draft called the LSTM
%  advantage over LightGBM "significant at 5% (p = 0.041)". That is FORBIDDEN
%  by HANDOFF 16A/A1: raw p = 0.0405 but Holm-corrected p = 0.1215 over the
%  21-test exploratory family, so it does NOT survive correction. Rewritten as
%  a lower point estimate that does not survive multiplicity correction.
%  LSTM, ReLU and Adam: LSTM already footnoted in ch1-2; ReLU/Adam are new.

\subsection{\lr{LSTM}}
\label{sec:3-7-2}

بازوی یادگیری عمیقِ این پژوهش یک شبکه‌ی حافظه‌ی کوتاه‌مدتِ بلند است. انگیزه‌ی انتخاب آن،
ساختار توالی‌وارِ سری قیمت است: قیمت هر روز نه یک بردار منفرد، بلکه حلقه‌ای از
زنجیره‌ای است که روزهای پیشین آن را شکل داده‌اند، و لایه‌ی بازگشتی سازوکاری صریح برای
مدل‌سازی همین وابستگی زمانی در اختیار می‌گذارد \cite{zhang_deep_2020,
ugurlu_electricity_2018} \lr{—} چیزی که نه مدل خطیِ \lr{LEAR} و نه مدل درختیِ
\lr{LightGBM} به‌صراحت ندارند و تنها از طریق ویژگی‌های وقفه‌ای به آن دسترسی می‌یابند.

\subsubsection*{معماری دوشاخه}

شبکه دو ورودی مجزا دارد که در میانه‌ی مسیر به هم می‌پیوندند.

\textbf{شاخه‌ی توالی.} به‌ازای هر وقفه‌ی قیمتیِ پیکربندی‌شده \lr{—} به‌ترتیب از کهنه به
تازه: \lr{D-7}، \lr{D-3}، \lr{D-2} و \lr{D-1} \lr{—} یک گام زمانی ساخته می‌شود و
محتوای هر گام، بردار ۲۴ساعته‌ی قیمت همان روزِ وقفه است. حاصل، دنباله‌ای چهارگامی با
بُعد ۲۴ در هر گام است که به لایه‌ی بازگشتی خورانده می‌شود.

انتخاب ستون‌ها صریحاً بر اساس نامِ ستون انجام می‌گیرد و هرگز بر پایه‌ی موقعیت آن؛ اگر
ستونی از قلم بیفتد، مدل با خطا متوقف می‌شود و به‌آرامی روی ستون اشتباه برازش
نمی‌شود. این تمایز اهمیت دارد، چون برازش بر ستون اشتباه هیچ نشانه‌ی بیرونی ندارد.

\textbf{شاخه‌ی ایستا.} تمام ستون‌های باقی‌مانده‌ی ماتریس ویژگی \lr{—} وقفه‌های متغیرهای
برون‌زا، پیش‌بینی‌های روزپیشِ خودِ روز هدف، و متغیرهای مجازیِ روز هفته \lr{—} به‌صورت
یک بردار مسطح وارد می‌شوند.

خروجی لایه‌ی بازگشتی با بردار ایستا الحاق می‌شود، از یک لایه‌ی چگال با تابع فعال‌سازِ
\lr{ReLU}\footnote{\lr{Rectified Linear Unit}} می‌گذرد، و در نهایت به لایه‌ی چگالِ
۲۴تایی می‌رسد که همان بردار قیمت ساعتیِ روز \lr{D+1} است. بهینه‌سازِ آموزش
\lr{Adam}\footnote{\lr{Adaptive Moment Estimation}} است.

نکته‌ی روش‌شناختیِ مهم در این معماری آن است که هر دو شاخه از \emph{همان} ماتریس
ویژگیِ مشترک تغذیه می‌شوند که سایر مدل‌ها مصرف می‌کنند. مسیر جداگانه‌ای برای ساخت
توالی وجود ندارد. اگر چنین مسیری ساخته می‌شد، سطح دومی برای نشت اطلاعات پدید می‌آمد
که آزمون \cref{sec:3-4} آن را پوشش نمی‌داد؛ در عوض، تنها یک منبع ویژگی وجود دارد و همان یک
منبع ممیزی‌شده است.

\subsubsection*{مقیاس‌بندی}

هر سه جزء \lr{—} دنباله‌ی ورودی، بردار ایستا و خودِ هدف \lr{—} استاندارد می‌شوند و
پیش‌بینی‌ها در انتها به مقیاس اصلی بازگردانده می‌شوند.

تصریح ضروری آن است که این مقیاس‌گرها \emph{تنها بر برشِ آموزشِ همان بازبرازش} برازش
می‌شوند، نه بر کل سری. برازش مقیاس‌گر بر تمام داده، میانگین و انحراف معیارِ دوره‌ی
آزمون را به مرحله‌ی آموزش نشت می‌داد؛ نشتی ظریف که هیچ خطایی تولید نمی‌کند و تنها
نتیجه را به‌آرامی خوش‌بینانه می‌سازد.

\subsubsection*{تابع زیان}

زیان آموزش، خطای مطلق میانگین است. این انتخاب دو کارکرد دارد: هم‌راستایی با معیار
سرشناسِ گزارش در فصل چهارم، و هم‌ترازی با تابع هدف \lr{LightGBM} (\cref{sec:3-7-1}). بدون
این هم‌ترازی، بخشی از اختلاف میان دو مدل به تفاوت در تابع زیان بازمی‌گشت و مقایسه،
به‌جای پاسخ به پرسش «کدام خانواده‌ی مدل بهتر است»، به مصنوعی از انتخاب زیان تبدیل
می‌شد.

\subsubsection*{آهنگ واسنجی و پاسداشت ترتیب زمانی}

بازآموزیِ کامل یک شبکه‌ی عصبی در هر یک از ۷۲۸ مبدأ، بر سخت‌افزار این پژوهش عملی نبود.
لذا بازبرازش کامل هر هفت روز انجام می‌شود و شبکه در فاصله‌ی میان دو بازبرازش بدون
تغییر به کار می‌رود. با این‌همه پیش‌بینی همچنان برای \emph{هر} مبدأ روزانه تولید
می‌شود، پس افق و شمار پیش‌بینی‌ها میان مدل‌ها یکسان می‌ماند. این انحراف از پروتکل
واسنجی روزانه، هم‌سنگ انحراف ثبت‌شده برای \lr{SARIMAX} (\cref{sec:3-6})، به‌صراحت در دفتر
تصمیمات ثبت شده است، نه به‌عنوان جزئیاتی پنهان در پیاده‌سازی.

افزون بر آن، پیاده‌سازی یک پاسبان صریح دارد: اگر پنجره‌ی آموزشِ درخواستی پایانی
\emph{پیش از} آخرین بازبرازش کامل داشته باشد، مدل با خطا متوقف می‌شود. چنین حالتی
به‌معنای پاسخ‌دادن به یک مبدأ گذشته با شبکه‌ای است که بر داده‌ی متأخرتر آموزش دیده
\lr{—} یعنی نشت اطلاعاتِ واقعی. چارچوب پیش‌رونده مبادی را یکنواخت و صعودی تولید
می‌کند، پس فعال‌شدن این پاسبان تنها می‌تواند نشانه‌ی اشکال در فراخوان باشد.

\subsubsection*{تنظیم ابرپارامترها}

همانند \lr{LightGBM}، ۵۰ آزمایش \lr{Optuna} با بذر نمونه‌بردار ۴۲، بر پنجره‌ی
اعتبارسنجیِ ۵ ژانویه‌ی ۲۰۱۵ تا ۳ ژانویه‌ی ۲۰۱۶ و پنجره‌ی برازشِ ۱۶ ژانویه‌ی ۲۰۱۲ تا ۴
ژانویه‌ی ۲۰۱۵ \lr{—} هر دو اکیداً پیش از دوره‌ی آزمون. بهترین پیکربندی: ۲۹ واحد پنهان،
۴۷ دوره‌ی آموزش، اندازه‌ی دسته‌ی ۳۲ و نرخ یادگیری \lr{0.00083}، با \lr{MAE} اعتبارسنجیِ
\lr{3.696}.

شایان توجه است که جست‌وجو شبکه‌ای کوچک‌تر از پیش‌فرض را برگزید \lr{—} ۲۹ واحد در
برابر ۶۴ \lr{—} همان جهت‌گیریِ منظم‌سازانه‌ای که در تنظیم \lr{LightGBM} نیز دیده شد،
و نشانه‌ی محدودبودن اطلاعات قابل استخراج از این سری در برابر ظرفیت مدل است.

\subsubsection*{جبرگرایی}

پیش از هر ساخت و برازش، بذر تصادفیِ سراسری ۴۲ تثبیت و جبرگراییِ عملیات فعال می‌شود؛
شمار نخ‌های پردازش نیز بر مقداری ثابت تنظیم می‌گردد تا هم نتیجه بازتولیدپذیر بماند و
هم اجراهای هم‌زمانِ چارچوب پیش‌رونده بی‌ثبات نشوند. اگر جبرگراییِ عملیات در محیطی
در دسترس نباشد، پیاده‌سازی هشدار صریح می‌دهد، به‌جای آنکه در سکوت تضمین
بازتولیدپذیری را از دست بدهد.

\subsubsection*{عملکرد، و آنچه درباره‌ی آن نباید ادعا شود}

بر دوره‌ی آزمون، \lr{MAE} ساعتی این مدل \lr{3.873} و \lr{rMAE} آن \lr{0.424} است
\lr{—} یعنی بهترین مدل منفرد در \cref{tab:hourly-results}. اما این عدد به‌تنهایی گمراه‌کننده است
و سه قید باید همراه آن بیاید.

نخست، آزمون دیبولد\lr{–}ماریانو با تصحیح \lr{HAC} نشان می‌دهد که برتری این مدل بر
\lr{LEAR-LASSO} \textbf{معنادار نیست}: \lr{p} برابر \lr{0.404}. این یک تساوی آماری
است، نه یک پیروزی، و در سرتاسر پایان‌نامه به همین صورت گزارش می‌شود.

برتری آن بر \lr{LightGBM} نیز، برخلاف آنچه نگاه نخست به \lr{p}-مقدار خام
(\lr{0.0405}) نشان می‌دهد، از تصحیح چندگانگی جان به در نمی‌برد: در خانواده‌ی اکتشافیِ
۲۱تایی که \cref{sec:3-5} آن را اعلام کرد، مقدار تصحیح‌شده با روش هولم\lr{–}بونفرونی برابر
\lr{0.1215} است. پس آنچه می‌توان گفت این است که \lr{LSTM} برآورد نقطه‌ای پایین‌تری
دارد، و این تفاوت پس از تصحیح برای چندگانگی دوام نمی‌آورد. خوانشِ صادقانه آن است که
\lr{LEAR-LASSO}، \lr{LightGBM} و \lr{LSTM} متقابلاً تفکیک‌ناپذیرند.

دوم، در آزمون برون‌توزیعی بر داده‌ی زنده‌ی ۲۰۲۶، همین مدل دومین افت شدید را پس از
\lr{LightGBM} نشان می‌دهد: \lr{rMAE} برابر \lr{1.520} در برابر \lr{1.087} برای
\lr{LEAR-LASSO}. یعنی مدلی که مزیت درون‌دوره‌اش از نظر آماری قابل تفکیک نبود، در
برابر تغییر شرایط شکننده‌تر نیز هست.

سوم، تحلیل \lr{SHAP} \cref{sec:4-6} بر یک برازش ایستا انجام شده و نه بر مدلی که در هر مبدأ
بازواسنجی می‌شود؛ لذا نمودارهای اهمیت ویژگی، تصویرِ همان مدلی که اعداد دقت را تولید
کرده است نیستند. آن برازش، مطابق تفکیکی که \cref{sec:3-5} میان سه مصنوع مدل قائل شد، پنجره‌ی
آموزشی اکیداً پیش از مرز دوره‌ی آزمون دارد و بنابراین هر روزِ توضیح‌داده‌شده برای آن
دیده‌نشده است؛ همزادی وفادار است، نه همان مدل. این قید در \cref{sec:4-6} نیز تصریح شده است.

جمع‌بندیِ صادقانه آن است که مدل عمیق در این پژوهش بهترین عدد میانگین را تولید کرد،
بدون آنکه برتری‌اش بر مدل خطیِ مرجع از حد نوسان آماری فراتر رود، و با هزینه‌ای آشکار
در تعمیم‌پذیری بیرون از دوره‌ی آموزش. تفصیل این بحث در بخش‌های ۴-۵ و ۵-۲ می‌آید.
